\documentclass[a4,11pt]{aleph-notas}
% Se puede ver la documentación aquí:
% https://github.com/alephsub0/LaTeX_aleph-notas

% -- Paquetes adicionales
\usepackage{enumitem}
\usepackage{aleph-comandos}
\usepackage{booktabs}


% -- Datos
\institucion{Facultad de Ciencias Exactas, Naturales y Ambientales}
\carrera{Catálogo STEM}
\asignatura{Métodos Numéricos}
\tema{Resumen no. 3: Sistemas de ecuaciones}
\autor{Mario Cueva - Andrés Merino}
\fecha{Periodo 2026-1}

\logouno[0.18\textwidth]{Logos/logoPUCE_04_ac}
\definecolor{colortext}{HTML}{1957d1}
\definecolor{colordef}{HTML}{1957d1}
\fuente{montserrat}


% -- Comandos adicionales
\setlist[enumerate]{label=\roman*.}
\newcommand{\vect}[1]{\mathbf{#1}}


\begin{document}

\encabezado

%%%%%%%%%%%%%%%%%%%%%%%%%%%%%%%%%%%%%%
\section{Resolución de sistemas de ecuaciones no lineales}

\begin{defi}[Sistema de ecuaciones no lineales]
    Un \textbf{sistema de ecuaciones no lineales} busca un vector
    \[
        x=
        \begin{pmatrix}
            x_1\\
            x_2\\
            \vdots\\
            x_n
        \end{pmatrix}
    \]
    que satisfaga
    \[
        F({x})=
        \begin{pmatrix}
            f_1(x_1,\ldots,x_n)\\
            f_2(x_1,\ldots,x_n)\\
            \vdots\\
            f_n(x_1,\ldots,x_n)
        \end{pmatrix}
        =
        {0}.
    \]
    El sistema es no lineal si al menos una de sus ecuaciones no es lineal en sus variables.
\end{defi}

\begin{ejem}
    El sistema
    \[
        \begin{cases}
            x^2+y^2-4=0,\\
            x-y=0
        \end{cases}
    \]
    es no lineal por la presencia de $x^2$ y $y^2$. En forma vectorial se escribe como
    \[
        F(x,y)=
        \begin{pmatrix}
            x^2+y^2-4\\
            x-y
        \end{pmatrix}=
        \begin{pmatrix}
            0\\
            0
        \end{pmatrix}.
    \]
\end{ejem}


\begin{defi}[Matriz jacobiana]
    Si
    \[
        F=(f_1,f_2,\ldots,f_m)
    \]
    depende de las variables $x_1,x_2,\ldots,x_n$, su \textbf{matriz jacobiana} es
    \[
        J_{F}({x})
        =
        \begin{pmatrix}
            \dfrac{\partial f_1}{\partial x_1} &
            \dfrac{\partial f_1}{\partial x_2} &
            \cdots &
            \dfrac{\partial f_1}{\partial x_n}\\[5mm]
            \dfrac{\partial f_2}{\partial x_1} &
            \dfrac{\partial f_2}{\partial x_2} &
            \cdots &
            \dfrac{\partial f_2}{\partial x_n}\\[2mm]
            \vdots & \vdots & \ddots & \vdots\\[2mm]
            \dfrac{\partial f_m}{\partial x_1} &
            \dfrac{\partial f_m}{\partial x_2} &
            \cdots &
            \dfrac{\partial f_m}{\partial x_n}
        \end{pmatrix}.
    \]
\end{defi}

\begin{ejem}
    Para
    \[
        F(x,y)=
        \begin{pmatrix}
            x^2+y^2-4\\
            x-y
        \end{pmatrix},
    \]
    se obtiene
    \[
        J_{{F}}(x,y)=
        \begin{pmatrix}
            2x & 2y\\
            1 & -1
        \end{pmatrix}.
    \]
\end{ejem}

\subsection{Método del punto fijo}

\begin{defi}[Punto fijo para sistemas]
    Un sistema no lineal puede transformarse en una ecuación vectorial de la forma
    \[
        G(x) = x.
    \]
    El método del punto fijo genera aproximaciones mediante
    \[
        {x}^{(k+1)}={G}\bigl({x}^{(k)}\bigr),
    \]
    a partir de una aproximación inicial ${x}^{(0)}$.
\end{defi}

\begin{teo}[Condición suficiente de convergencia]
    Sea $D$ una región que contiene al punto fijo y supongamos que ${G}(D)\subseteq D$. Si existe una norma matricial y una constante $0<q<1$ tales que
    \[
        \left\|J_{{G}}({x})\right\|\leq q
    \]
    para todo ${x}\in D$, entonces la iteración de punto fijo converge en $D$ hacia un único punto fijo.
\end{teo}

\begin{advertencia}
    Escribir el sistema como ${x}={G}({x})$ no basta para asegurar convergencia. Distintas transformaciones equivalentes pueden producir iteraciones que convergen, oscilan o divergen.
\end{advertencia}

\begin{ejem}
    Consideremos
    \[
        \begin{cases}
            3x-y-1=0,\\
            4y-x^2-1=0.
        \end{cases}
    \]
    Una forma de punto fijo es
    \[
        {G}(x,y)=
        \begin{pmatrix}
            \dfrac{1+y}{3}\\[2mm]
            \dfrac{1+x^2}{4}
        \end{pmatrix}.
    \]
    Desde ${x}^{(0)}=(0,0)$ se obtiene
    \[
        {x}^{(1)}=
        \begin{pmatrix}
            1/3\\
            1/4
        \end{pmatrix},
        \qquad
        {x}^{(2)}=
        \begin{pmatrix}
            5/12\\
            5/18
        \end{pmatrix}.
    \]
    El proceso continúa hasta que dos aproximaciones consecutivas difieran menos que la tolerancia elegida.
\end{ejem}

\subsection{Método de Newton-Raphson}

\begin{defi}[Newton-Raphson para sistemas]
    Para resolver
    \[
        {F}({x})={0},
    \]
    el método de Newton-Raphson usa, si el jacobiano es invertible:
    \[
        {x}^{(k+1)}
        =
        {x}^{(k)}
        -
        J_{{F}}\bigl({x}^{(k)}\bigr)^{-1}
        {F}\bigl({x}^{(k)}\bigr).
    \]
\end{defi}

\begin{ejem}
    Para
    \[
        {F}(x,y)=
        \begin{pmatrix}
            x^2+y^2-4\\
            x-y
        \end{pmatrix}
    \]
    y ${x}^{(0)}=(1,1)$ se tiene
    \[
        {F}(1,1)=
        \begin{pmatrix}
            -2\\
            0
        \end{pmatrix},
        \qquad
        J_{{F}}(1,1)=
        \begin{pmatrix}
            2 & 2\\
            1 & -1
        \end{pmatrix}.
    \]
        Además,
    \[
        J_{{F}}(1,1)^{-1}
        =
        \frac{1}{-4}
        \begin{pmatrix}
            -1 & -2\\
            -1 & 2
        \end{pmatrix}
        =
        \begin{pmatrix}
            \frac{1}{4} & \frac{1}{2}\\
            \frac{1}{4} & -\frac{1}{2}
        \end{pmatrix}.
    \]
    Por el método de Newton,
    \[
        {x}^{(1)}
        =
        {x}^{(0)}-J_{{F}}(1,1)^{-1}{F}(1,1).
    \]
    Entonces,
    \[
        J_{{F}}(1,1)^{-1}{F}(1,1)
        =
        \begin{pmatrix}
            \frac{1}{4} & \frac{1}{2}\\
            \frac{1}{4} & -\frac{1}{2}
        \end{pmatrix}
        \begin{pmatrix}
            -2\\
            0
        \end{pmatrix}
        =
        \begin{pmatrix}
            -\frac{1}{2}\\
            -\frac{1}{2}
        \end{pmatrix}.
    \]
    Por tanto,
    \[
        {x}^{(1)}
        =
        \begin{pmatrix}
            1\\
            1
        \end{pmatrix}
        -
        \begin{pmatrix}
            -\frac{1}{2}\\
            -\frac{1}{2}
        \end{pmatrix}
        =
        \begin{pmatrix}
            \frac{3}{2}\\
            \frac{3}{2}
        \end{pmatrix}.
    \]
    Así, la primera iteración del método de Newton da
    \[
        {x}^{(1)}=\left(\frac{3}{2},\frac{3}{2}\right).
    \]
\end{ejem}

%%%%%%%%%%%%%%%%%%%%%%%%%%%%%%%%%%%%%%
\section{Resolución de sistemas de ecuaciones lineales}

\subsection{Método de Gauss}

\begin{defi}[Sistema lineal]
    Un sistema lineal de $n$ ecuaciones con $n$ incógnitas puede escribirse como
    \[
        A{x}={b},
    \]
    donde $A$ es la matriz de coeficientes, ${x}$ es el vector de incógnitas y ${b}$ es el vector de  resultados.
\end{defi}

\begin{defi}[Eliminación de Gauss]
    El \textbf{método de Gauss} transforma la matriz aumentada
    \[
        [A\mid {b}]
    \]
    en una matriz triangular superior mediante operaciones elementales por filas. Luego se obtiene la solución por sustitución regresiva.
\end{defi}

\begin{ejem}
    Para resolver
    \[
        \begin{cases}
            2x+y-z=8,\\
            -3x-y+2z=-11,\\
            -2x+y+2z=-3,
        \end{cases}
    \]
    se trabaja con
    \[
        \left[
        \begin{array}{ccc|c}
            2 & 1 & -1 & 8\\
            -3 & -1 & 2 & -11\\
            -2 & 1 & 2 & -3
        \end{array}
        \right].
    \]
    Al eliminar los elementos bajo los pivotes se obtiene una forma triangular equivalente:
    \[
        \left[
        \begin{array}{ccc|c}
            2 & 1 & -1 & 8\\
            0 & 1/2 & 1/2 & 1\\
            0 & 0 & -1 & 1
        \end{array}
        \right].
    \]
    La sustitución regresiva produce
    \[
        z=-1,\qquad y=3,\qquad x=2.
    \]
\end{ejem}

\subsection{Condicionamiento y pivoteo}

\begin{defi}[Sistema mal condicionado]
    Un sistema está \textbf{mal condicionado} cuando pequeños cambios en los datos pueden causar cambios grandes en la solución. En ese caso, los errores de redondeo o de medición pueden amplificarse durante el cálculo.
\end{defi}

\begin{ejem}
    El sistema
    \[
        \begin{cases}
            x+y=2,\\
            x+1.0001y=2.0001
        \end{cases}
    \]
    tiene solución $(x,y)=(1,1)$. Si solo se cambia el segundo término independiente a $2.0002$, se obtiene
    \[
        \begin{cases}
            x+y=2,\\
            x+1.0001y=2.0002,
        \end{cases}
    \]
    cuya solución es $(x,y)=(0,2)$. Una perturbación pequeña en los datos produjo un cambio grande en la solución.
\end{ejem}

\begin{defi}[Número de condición]
    Para una matriz invertible $A$, el \textbf{número de condición} asociado a una norma se define por
    \[
        \kappa(A)=\|A\|\|A^{-1}\|.
    \]
    Si $\kappa(A)$ es cercano a $1$, el problema suele ser estable frente a perturbaciones. Si $\kappa(A)$ es grande, la solución puede ser muy sensible a errores en los datos.
\end{defi}

\begin{advertencia}
    El número de condición describe la sensibilidad del problema, no solo la calidad de un algoritmo. Un algoritmo estable no elimina el efecto de trabajar con una matriz muy mal condicionada.
\end{advertencia}

\begin{defi}[Pivoteo parcial]
    En la eliminación de Gauss con \textbf{pivoteo parcial}, antes de eliminar en la columna $k$ se intercambia la fila actual con una fila inferior cuyo elemento en esa columna tenga mayor valor absoluto. Así se usa como pivote un elemento más conveniente.
\end{defi}

\begin{ejem}
    En la matriz
    \[
        \begin{pmatrix}
            0.0001 & 1\\
            1 & 1
        \end{pmatrix},
    \]
    usar $0.0001$ como primer pivote introduce multiplicadores grandes. Con pivoteo parcial se intercambian las filas y se usa el pivote $1$, lo cual reduce la amplificación de errores de redondeo.
\end{ejem}

\subsection{Linealización}

\begin{defi}[Linealización de una ecuación]
    Si $f$ es derivable cerca de $x_0$, entonces
    \[
        f(x)\approx f(x_0)+f'(x_0)(x-x_0).
    \]
    Al reemplazar una ecuación no lineal $f(x)=0$ por esta aproximación lineal se obtiene
    \[
        f(x_0)+f'(x_0)(x-x_0)=0.
    \]
    Resolver esta ecuación para $x$ produce precisamente un paso del método de Newton.
\end{defi}

\begin{ejem}
    Para $f(x)=x^2-2$ alrededor de $x_0=1.4$,
    \[
        f(1.4)=-0.04,
        \qquad
        f'(1.4)=2.8.
    \]
    La ecuación linealizada es
    \[
        -0.04+2.8(x-1.4)=0,
    \]
    y su solución es
    \[
        x\approx 1.4143.
    \]
\end{ejem}

\begin{defi}[Linealización de un sistema]
    Para un sistema ${F}({x})={0}$ y un punto ${x}^{(0)}$, la aproximación lineal es
    \[
        {F}({x})
        \approx
        {F}\bigl({x}^{(0)}\bigr)
        +
        J_{{F}}\bigl({x}^{(0)}\bigr)
        \bigl({x}-{x}^{(0)}\bigr).
    \]
    Al igualar esta expresión a ${0}$ se obtiene un sistema lineal para la corrección ${x}-{x}^{(0)}$.
\end{defi}

\begin{ejem}
    Para
    \[
        \begin{cases}
            x^2+y^2-4=0,\\
            x-y=0
        \end{cases}
    \]
    alrededor de $(1,1)$, la linealización queda
    \[
        \begin{cases}
            -2+2(x-1)+2(y-1)=0,\\
            (x-1)-(y-1)=0.
        \end{cases}
    \]
    Al resolver este sistema lineal se obtiene $(x,y)=(1.5,1.5)$, la misma primera aproximación que entrega Newton-Raphson.
\end{ejem}



%%%%%%%%%%%%%%%%%%%%%%%%%%%%%%%%%%%%%%
\section{Factorización LU e inversión de matrices}

\subsection{Factorización LU}

\begin{defi}[Factorización LU]
    Una \textbf{factorización LU} escribe una matriz cuadrada como
    \[
        A=LU,
    \]
    donde $L$ es triangular inferior y $U$ es triangular superior. Cuando no se requieren intercambios de filas, los multiplicadores usados en eliminación de Gauss forman la matriz $L$ y la matriz triangular resultante es $U$.
\end{defi}

\begin{defi}[Resolución mediante LU]
    Si $A=LU$, resolver
    \[
        A{x}={b}
    \]
    equivale a resolver dos sistemas triangulares:
    \[
        L{y}={b},
        \qquad
        U{x}={y}.
    \]
    Primero se aplica sustitución progresiva y luego sustitución regresiva.
\end{defi}

\begin{advertencia}
    Si la eliminación necesita intercambios de filas, la forma usual es
    \[
        PA=LU,
    \]
    donde $P$ registra las permutaciones realizadas.
\end{advertencia}

\begin{ejem}
    La matriz
    \[
        A=
        \begin{pmatrix}
            2 & 1 & 1\\
            4 & -6 & 0\\
            -2 & 7 & 2
        \end{pmatrix}
    \]
    admite la factorización
    \[
        A=
        \begin{pmatrix}
            1 & 0 & 0\\
            2 & 1 & 0\\
            -1 & -1 & 1
        \end{pmatrix}
        \begin{pmatrix}
            2 & 1 & 1\\
            0 & -8 & -2\\
            0 & 0 & 1
        \end{pmatrix}.
    \]
    En este caso, los valores $2$, $-1$ y $-1$ que aparecen bajo la diagonal de $L$ corresponden a los multiplicadores de la eliminación.
\end{ejem}

\subsection{Inversión de matrices}

\begin{defi}[Matriz inversa]
    Una matriz cuadrada $A$ es invertible si existe una matriz $A^{-1}$ tal que
    \[
        AA^{-1}=A^{-1}A=I.
    \]
    Para calcular $A^{-1}$ se puede resolver
    \[
        A{x}_j={e}_j,
        \qquad j=1,\ldots,n,
    \]
    donde ${e}_j$ son las columnas de la identidad. Las soluciones ${x}_j$ forman las columnas de $A^{-1}$.
\end{defi}

\begin{advertencia}
    Para resolver un único sistema $A{x}={b}$ suele ser preferible resolver el sistema directamente o usar una factorización de $A$, en lugar de calcular la inversa completa.
\end{advertencia}

\begin{ejem}
    Para
    \[
        A=
        \begin{pmatrix}
            2 & 1\\
            5 & 3
        \end{pmatrix},
    \]
    se tiene $\det(A)=1$, por lo que
    \[
        A^{-1}=
        \begin{pmatrix}
            3 & -1\\
            -5 & 2
        \end{pmatrix}.
    \]
    En efecto,
    \[
        \begin{pmatrix}
            2 & 1\\
            5 & 3
        \end{pmatrix}
        \begin{pmatrix}
            3 & -1\\
            -5 & 2
        \end{pmatrix}
        =
        \begin{pmatrix}
            1 & 0\\
            0 & 1
        \end{pmatrix}.
    \]
\end{ejem}


\end{document}
